Machine learning models are now widely used for molecular property prediction (MPP), offering a data-driven alternative to quantum chemical simulation and experimental screening in areas such as drug discovery, materials science, and photophysical materials design~\cite{butler2018machine, schneider2020rethinking, jha2018elemnet}. Deep learning approaches that learn structure–property relationships directly from molecular graphs have achieved state-of-the-art performance across a wide array of MPP benchmarks~\cite{ying2021transformers, hu2020open, gasteiger2020directional}.

Despite these successes, the practical utility of deep models is frequently constrained by data availability. Experimental research groups often operate in a ``low-data'' regime, generating task-specific datasets ranging from tens to a few thousand molecules. A variety of deep learning architectures have been applied to MPP tasks, including sequence-based models operating on string-based representations and graph-based models operating directly on molecular structure. Among these, graph neural networks (GNNs) have emerged as a dominant paradigm because they incorporate molecular topology as an explicit inductive bias. However, in low-data settings, end-to-end GNNs often exhibit unstable training dynamics, high sensitivity to random initialization, and poor generalization under distribution shift~\cite{chuang2020debiased, grinsztajn2022tree}. This raises critical questions about the efficacy of GNN-based deep learning relative to simpler baselines when labeled data is scarce.

These challenges are particularly prevalent in the prediction of optical properties, such as absorption and emission wavelengths, which are vital for organic electronics, photonic materials, and fluorescent probes~\cite{facchetti2011pi}. Unlike ground-state properties (e.g., solubility or toxicity), optical responses are governed by excited-state electronic structure, which depends sensitively on $\pi$-conjugation length, solvent environment, and conformational flexibility~\cite{lakowicz2006principles, martin2020electronic, reichardt2011solvents}. Small structural modifications can yield large, non-intuitive spectral shifts, creating ``activity cliffs'' that are difficult for smooth, continuous models to capture without dense sampling~\cite{van2022exposing}.

While quantum chemical methods such as Time-Dependent Density Functional Theory (TD-DFT) provide physically grounded estimates of excited-state properties, they scale poorly with both molecule and dataset size, requiring substantial computational resources~\cite{dreuw2005single, casida2009time}. Tannir \textit{et al.} demonstrate that descriptor-based gradient-boosted ensembles can achieve strong performance on curated optical datasets when augmented with handcrafted TD-DFT–derived features~\cite{tannir2024predicting}. However, such approaches depend on computing quantum-chemical descriptors for each molecule at inference time, limiting their applicability in large-scale screening. To retain the informational benefits of quantum calculations without incurring inference-time costs, recent work by Greenman \textit{et al.} has explored multi-fidelity strategies that train machine learning surrogates jointly on experimental and computational datasets to avoid the need for DFT calculations at inference~\cite{greenman2022multi, ramakrishnan2015big}.

Conversely, in purely experimental low-data regimes, classical models based on fixed molecular descriptors (e.g., Morgan fingerprints) often outperform deep learning approaches due to their strong inductive biases and statistical efficiency~\cite{rogers2010extended, wu2018moleculenet, xia2023understanding}. However, these approaches may fail to capture the higher-order structural effects and long-range dependencies relevant to optical response.

Transfer learning has been proposed as a remedy for the lack of labeled data, but its effectiveness for molecular regression tasks is highly regime-dependent. Recent studies indicate that without strict chemical alignment between pretraining and downstream tasks, transfer learning can yield negligible or even negative transfer, particularly in regression tasks involving distinct physical mechanisms~\cite{hu2019strategies, sun2022does, buterez2024transfer}.

To balance the expressive power of GNNs with the statistical stability of classical regression, researchers are increasingly adopting hybrid modelling strategies. For instance, Deng \textit{et al.} introduced XGraphBoost, a framework that utilizes GNN-learned representations as input features for Gradient-Boosted Decision Trees (GBDTs). Their comprehensive benchmark evaluation identified the combination of eXtreme Gradient Boosting (XGB) with a Directed Message Passing Neural Network (D-MPNN) as the top-performing hybrid strategy~\cite{deng2021xgraphboost,chen2016xgboost,yang2019analyzing}. While this work establishes the promise of hybrid architectures, it does not address the regimes most relevant to experimental photophysics, where datasets are small, domain-specific, and often subject to distribution shift. Moreover, prior hybrid studies have largely focused on aggregate benchmark performance, without examining how predictive accuracy, variance, and learned representations evolve as a function of dataset size.

Hybrid models explicitly decouple representation learning from downstream regression, making the structure of the learned embedding space a central object of study rather than a hidden intermediate. In low-data regimes, where models are prone to memorizing noise rather than learning generalizable features, it is imperative to ensure that learned representations capture chemically meaningful structure rather than overfitting to the limited training data. 

In this work, we systematically evaluate a suite of classical, D-MPNN, and hybrid modelling strategies for optical property prediction across multiple datasets, data regimes, and train–test splits. We compare end-to-end D-MPNNs with hybrid models that use D-MPNN embeddings to train lightweight tree-based heads. We demonstrate that hybrid models broadly outperform end-to-end D-MPNNs on peak absorption wavelength prediction in small-data regimes. Through controlled transfer learning experiments, we show that the benefits of pretraining are strongly regime- and coverage-dependent, yielding substantial gains on chemically aligned downstream datasets but limited or negative returns under significant distribution shift. Furthermore, using a nested experimental design, we decompose predictive variability into contributions arising from representation learning and downstream prediction; we find that while predictive instability is dominated by representation learning across regimes, the downstream head contributes a substantial fraction in data-limited settings. Finally, we analyze the structure of learned molecular embeddings and demonstrate that these representations retain chemically meaningful organization aligned with known determinants of optical response, providing practical guidance for modelling optical properties in small and specialized chemical datasets.
